%% Weak-Memory Local Reasoning, Revisited.
%% Build: latexmk -pdf main   (or: tectonic main.tex)
\documentclass[11pt]{article}
\usepackage[T1]{fontenc}
\usepackage[margin=1in]{geometry}
\usepackage{microtype}
\usepackage{booktabs}
\usepackage{tabularx}
\usepackage{enumitem}
\usepackage{mathpartir}
\usepackage{xcolor}
\usepackage{placeins}
\usepackage[bookmarks=true,colorlinks=true,linkcolor=black,citecolor=black,urlcolor=blue]{hyperref}

%% The dissertation's own notation, unchanged.
%!TEX root = dissertation.tex

\usepackage{amsmath,amssymb,amsthm}
\usepackage{graphicx}
\usepackage{stmaryrd}
\usepackage{nicefrac}

\theoremstyle{plain}
\newtheorem{lemma}{Lemma}
\newtheorem{theorem}{Theorem}
\newtheorem{corollary}{Corollary}
\newtheorem{proposition}{Proposition}
\newtheorem{conjecture}{Conjecture}

\theoremstyle{definition}
\newtheorem{definition}{Definition}

\theoremstyle{remark}
\newtheorem{note}{Note}
\newtheorem*{remark}{Remark}


\newcommand{\bnfbar}{\ensuremath{~|~}}
\newcommand{\bnfdef}{\ensuremath{~::=~}}

%% set notation
\newcommand{\set}[1]{\left\lbrace #1 \right\rbrace}
\newcommand{\setof}[2]{\set{#1 \mid #2}}
\newcommand{\bigset}[1]{\left\lbrace #1 \right\rbrace}
\newcommand{\bigsetof}[2]{\bigset{#1 \mid #2}}

\DeclareMathOperator{\funcompat}{\smile}

\newcommand{\mtab}{\ensuremath{\;\;\;}}

\newcommand{\powerset}[1]{\ensuremath{\mathcal{P}(#1)}}
\newcommand{\zorch}{\ensuremath{\lightning}}

\DeclareMathOperator{\opdom}{\mathrm{dom}}
\newcommand{\dom}[1]{\ensuremath{\opdom(#1)}}

\newcommand{\bvt}{\ensuremath{\mathbf{t}}}
\newcommand{\bvf}{\ensuremath{\mathbf{f}}}
\newcommand{\llen}[1]{\ensuremath{|#1|}}
\newcommand{\llit}[1]{\ensuremath{\left[#1\right]}}
\newcommand{\lnil}[0]{\ensuremath{\varepsilon}}
\newcommand{\arnil}[0]{\ensuremath{\mathcal{E}}}
\newcommand{\elems}[1]{\ensuremath{|\!|#1|\!|}}
\newcommand{\funof}[1]{\ensuremath{\overline{#1}}}
\DeclareMathOperator{\lapp}{+\!\!\!+}
\DeclareMathOperator{\lcons}{:\!:}
\newcommand{\lhead}[1]{\ensuremath{\mathit{hd}(#1)}}
\newcommand{\ltail}[1]{\ensuremath{\mathit{tl}(#1)}}
\newcommand{\nth}[2]{\ensuremath{#2 ! #1}}
\newcommand{\listof}[1]{\ensuremath{#1~\mathsf{list}}}

\newcommand{\relset}[1]{\ensuremath{|#1|}}
\newcommand{\ext}[1]{\ensuremath{\hat{#1}}}
\newcommand{\restrict}[2]{\ensuremath{#1 |_{#2}}}
\newcommand{\proj}[2]{\ensuremath{#1 |_{#2}}}

\newcommand{\reclit}[1]{\ensuremath{\set{#1}}}

\DeclareMathOperator{\tfun}{\rightarrow}
\DeclareMathOperator{\pfun}{\rightharpoonup}
\DeclareMathOperator{\fpfun}{\pfun_{\mathrm{fin}}}

\DeclareMathOperator{\st}{\,:\,}

\newcommand{\ptup}[2]{\ensuremath{#1 \shortleftarrow #2}}
\newcommand{\funup}[2]{\ensuremath{#1\! \left[ #2 \right]}}
\newcommand{\fundel}[2]{\ensuremath{#1 \setminus#2 }}
\newcommand{\recup}[3]{\funup{#1}{\ptup{#2}{#3}}}
\newcommand{\ptapp}[2]{\ensuremath{#1 \lapp #2}}

\DeclareMathOperator{\conj}{\,\wedge\,}
\DeclareMathOperator{\disj}{\,\vee\,}
\DeclareMathOperator{\onlyif}{\,\Rightarrow\,}
%\DeclareMathOperator{\ifandonlyif}{\,\Longleftrightarrow\,}
\renewcommand{\iff}{\ensuremath{\Leftrightarrow}}

\newcommand{\fresh}[1]{\ensuremath{\mathrm{fresh}(#1)}}
\newcommand{\subst}[2]{\ensuremath{\left[#1 / #2\right]}}
\newcommand{\unsubst}[2]{\ensuremath{\left[#1 \backslash #2\right]}}

\newcommand{\eqclass}[1]{\ensuremath{\left[ #1 \right]}}
\DeclareMathOperator{\downop}{\downarrow}
\newcommand{\downo}[2]{\ensuremath{\downop_{#1}\!#2}}
\newcommand{\down}[1]{\ensuremath{\downop #1}}

%% Calculational proofs

\newcommand{\Calc}[1]{\begin{description}
                \item \begin{tabbing}\qquad\=\quad\=\kill
                \> \> #1\end{tabbing}
                \end{description}}
\newcommand{\Sketch}[1]{\begin{description}
                \item \begin{tabbing}\qquad\=\quad\=\kill
                #1\end{tabbing}
                %\> \> #1\end{tabbing}
                \end{description}}                

\newcommand{\aconn}[1]{\{#1\}\\ }
\newcommand{\aconnl}[1]{\{#1\} }
\newcommand{\aconnt}[1]{\{#1\} $\therefore$\\ }
\newcommand{\aconni}[1]{\> \{#1\}\\ }
\newcommand{\acomm}[1]{\> #1 \\}


\newcommand{\conn}[2]{\\*$#1 $\> \{#2\}\\\> \>}
\newcommand{\Conn}[1]{\conn{=}{#1}}

\newcommand{\ident}[1]{\ensuremath{\mathit{Id}_{#1}}}
\DeclareMathOperator{\rcomp}{\circ}
\newcommand{\relexp}[2]{\ensuremath{\stackrel{#2}{#1}}}
\newcommand{\tcl}[1]{\overset{+}{#1}}
\newcommand{\rtcl}[1]{\overset{\ast}{#1}}

\newcommand{\image}[2]{\ensuremath{#1\left[#2\right]}}
\newcommand{\invimage}[2]{\ensuremath{#1^{-1}\left[#2\right]}}

\newcommand{\lift}[1]{\ensuremath{\mathit{lift}_{#1}}}
\renewcommand{\merge}{\ensuremath{\uplus}}

% \newcommand{\override}{\reflectbox{\rotatebox[origin=c]{180}{\ensuremath{\,\oslash\,}}}}
\newcommand{\override}{\ensuremath{\backslash\!\!\backslash}}

\newcommand{\sentails}{\ensuremath{\,\models\,}}
\newcommand{\satisfies}{\ensuremath{\,\models\,}}
\newcommand{\sequiv}{\ensuremath{\,\equiv\,}}

\newcommand{\fixme}[1]{\todo{FIXME: #1}}

\DeclareMathOperator{\eqdef}{\;=_{\mathit{df}}\;}
\DeclareMathOperator{\iffdef}{\;\equiv_{\mathit{df}}\;}

\newcommand{\card}[1]{\left|#1\right|}
\newcommand{\cmpl}[1]{\ensuremath{\overline{#1}}}

\newcommand{\refines}{\ensuremath{\sqsubseteq}}
\newcommand{\range}[1]{\ensuremath{\mathit{range}(#1)}}

\newcommand{\fun}[1]{\ensuremath{\lambda#1\,.\,}}
%!TEX root = dissertation.tex

\usepackage{stackrel}

%% types

\newcommand{\setlocations}{\ensuremath{\mathbb{L}}}
\newcommand{\setvalues}{\ensuremath{\mathbb{V}}}
\newcommand{\setintegers}{\ensuremath{\mathbb{Z}}}
\newcommand{\setnaturals}{\ensuremath{\mathbb{N}}}
\newcommand{\setpositives}{\ensuremath{\mathbb{N}^+}}
\newcommand{\setidentifiers}{\ensuremath{\mathbb{I}}}
\newcommand{\setbooleans}{\ensuremath{\mathbb{B}}}
\newcommand{\setprocessors}{\ensuremath{\mathbb{P}}}

%% abbreviations

\newcommand{\setstacks}{\ensuremath{\mathbf{Stack}}}
\newcommand{\setheaps}{\ensuremath{\mathbf{Heap}}}
\newcommand{\setbuffers}{\ensuremath{\mathbf{Buf}}}
\newcommand{\setshares}{\ensuremath{\mathbf{Share}}}
\newcommand{\setperms}{\ensuremath{\mathbf{Perm}}}
\newcommand{\setastates}{\ensuremath{\mathit{\Sigma}}}
\newcommand{\setstates}{\ensuremath{\mathbf{State}}}
\newcommand{\setmodels}{\ensuremath{\mathbf{Model}}}
\newcommand{\setmemorysystems}{\ensuremath{\mathbf{Mem}}}
\newcommand{\setpstates}{\ensuremath{\Pi}}
\newcommand{\setmpstates}{\ensuremath{\mathbf{MPState}}}
\newcommand{\setresolvers}{\ensuremath{\mathbf{Resolver}}}
\newcommand{\setconfigurations}{\ensuremath{\mathbf{Config}}}
\newcommand{\setmconfigurations}{\ensuremath{\mathbf{MConfig}}}
\newcommand{\setassertions}{\ensuremath{\mathbf{Asrt}}}

%% operators 

\DeclareMathOperator{\op}{\mathit{op}}
\DeclareMathOperator{\opskip}{\mathsf{skip}}
\DeclareMathOperator{\opfence}{\mathsf{fence}}
\DeclareMathOperator{\opassume}{\mathsf{assume}}
\DeclareMathOperator{\opassert}{\mathsf{assert}}
\DeclareMathOperator{\opassign}{:\!=}
\DeclareMathOperator{\opequals}{=}
\DeclareMathOperator{\opseq}{;}
\DeclareMathOperator{\opchoice}{+}
\DeclareMathOperator{\oppar}{|\!|}
\DeclareMathOperator{\oploop}{\ast}
\DeclareMathOperator{\oplock}{\mathsf{lock}}
\DeclareMathOperator{\opunlock}{\mathsf{unlock}}
\DeclareMathOperator{\opcons}{\mathsf{cons}}
\DeclareMathOperator{\opnew}{\mathsf{new}}
\DeclareMathOperator{\opalloc}{\mathsf{alloc}}
\DeclareMathOperator{\opdispose}{\mathsf{free}}
\DeclareMathOperator{\opif}{\mathsf{if\,}}
\DeclareMathOperator{\opthen}{\mathsf{\,then\,}}
\DeclareMathOperator{\opelse}{\mathsf{\,else}}
\DeclareMathOperator{\opwhile}{\mathsf{while\,}}
\DeclareMathOperator{\opdo}{\mathsf{\,do\,}}
\DeclareMathOperator{\oplocal}{\mathsf{local\,}}
\DeclareMathOperator{\opin}{\mathsf{\,in\,}}
\DeclareMathOperator{\opresource}{\mathsf{resource\,}}
\DeclareMathOperator{\opwith}{\mathsf{with\,}}
\DeclareMathOperator{\opwhen}{\mathsf{\,when\,}}
\DeclareMathOperator{\nil}{\emptyset}
\DeclareMathOperator{\compat}{\smile}


%% primitive commands

\newcommand{\cskip}[0]{\ensuremath{\opskip}}
\newcommand{\cfence}[0]{\ensuremath{\opfence}}
\newcommand{\clock}[0]{\ensuremath{\oplock}}
\newcommand{\cunlock}[0]{\ensuremath{\opunlock}}
\newcommand{\cassume}[1]{\ensuremath{\opassume(#1)}}
\newcommand{\cassert}[1]{\ensuremath{\opassert(#1)}}
\newcommand{\cassign}[2]{\ensuremath{#1 \opassign #2}}
\newcommand{\cload}[2]{\ensuremath{#1 \opassign \left[ #2 \right]}}
\newcommand{\cstore}[2]{\ensuremath{\left[ #1 \right] \opassign #2}}
\newcommand{\cnew}[2]{\ensuremath{#1 \opassign \opnew(#2)}}
\newcommand{\calloc}[1]{\ensuremath{#1 \opassign \opalloc}}
\newcommand{\ccons}[3]{\ensuremath{#1 \opassign \opcons(#2, #3)}}
\newcommand{\cfree}[1]{\ensuremath{\opdispose(#1)}}

%% commands

\newcommand{\cseq}[2]{\ensuremath{#1 \opseq #2}}
\newcommand{\cchoice}[2]{\ensuremath{#1 \opchoice #2}}
\newcommand{\cpar}[2]{\ensuremath{#1 \oppar #2}}
\newcommand{\cifthenelse}[3]{\ensuremath{\opif #1 \opthen #2 \opelse #3}}
\newcommand{\cifthen}[2]{\ensuremath{\opif #1 \opthen #2}}
\newcommand{\cwhile}[2]{\ensuremath{\opwhile #1 \opdo #2}}
\newcommand{\cloop}[1]{\ensuremath{#1^{\oploop}}}
\newcommand{\clocal}[3]{\ensuremath{\oplocal #1 \opassign #2 \opin #3}}
\newcommand{\cres}[2]{\ensuremath{\opresource #1 \opin #2}}
\newcommand{\cwith}[2]{\ensuremath{\left\langle #2\right\rangle_{#1}}}
% \newcommand{\cwith}[2]{\ensuremath{\opwith #1 \opdo #2}}
\newcommand{\catomic}[1]{\ensuremath{\left\langle #1 \right\rangle}}
\newcommand{\ccas}[2]{\ensuremath{\mathsf{cas}_{#1}(#2)}}

\newcommand{\exprs}{\ensuremath{\mathbf{Expr}}}
\newcommand{\fracs}{\ensuremath{\mathbf{Frac}}}
\newcommand{\bexprs}{\ensuremath{\mathbf{BExpr}}}
\newcommand{\pcomms}{\ensuremath{\mathbf{PComm}}}
\newcommand{\comms}{\ensuremath{\mathbf{Comm}}}
\newcommand{\asserts}{\ensuremath{\mathbf{Assert}}}
\newcommand{\progs}{\ensuremath{\mathbf{Prog}}}

\newcommand{\dnexpr}[1]{\ensuremath{\mathcal{E}[\![#1]\!]}}
\newcommand{\dnshare}[1]{\ensuremath{\mathcal{S}[\![#1]\!]}}
\newcommand{\dnbexpr}[1]{\ensuremath{\mathcal{B}[\![#1]\!]}}
\newcommand{\dnpcomm}[1]{\ensuremath{\mathcal{P}[\![#1]\!]}}
\newcommand{\dna}[1]{\ensuremath{\mathcal{A}[\![#1]\!]}}
\newcommand{\dnfrac}[1]{\ensuremath{\mathcal{F}[\![#1]\!]}}

\newcommand{\sflush}[1]{\ensuremath{\mathit{flush}(#1)}}
\newcommand{\speek}[2]{\ensuremath{\mathit{peek}_{#1}(#2)}}
\newcommand{\sval}[1]{\ensuremath{\mathit{val}(#1)}}
\newcommand{\spush}[3]{\ensuremath{\mathit{push}_{#1}^{#2}(#3)}}
\newcommand{\alloc}[1]{\ensuremath{\mathit{alloc}(#1)}}
\newcommand{\allocat}[2]{\ensuremath{\mathit{alloc}_{#1}(#2)}}
\newcommand{\owned}[1]{\ensuremath{\mathit{owned}(#1)}}
\newcommand{\full}[1]{\ensuremath{\mathit{full}(#1)}}

\newcommand{\step}{\ensuremath{\rightarrow}}
\newcommand{\rtstep}{\ensuremath{\stackrel{\ast}{\rightarrow}}}
\newcommand{\nostep}{\ensuremath{\nrightarrow}}
\newcommand{\pstep}[1]{\ensuremath{\underset{#1}{\rightarrow}}}
\newcommand{\taustep}{\pstep{\tau}}
\newcommand{\modstep}{\pstep{\widehat{\tau}}}

\newcommand{\fv}[1]{\ensuremath{\mathrm{fv}(#1)}}
\renewcommand{\mod}[1]{\ensuremath{\mathrm{mod}(#1)}}

\newcommand{\bexpt}{\ensuremath{\mathbf{true}}}
\newcommand{\bexpf}{\ensuremath{\mathbf{false}}}


\newcommand{\spec}[4]{\ensuremath{#1\,\vdash\,\set{#2}\;#3\;\set{#4}}}
\newcommand{\truespec}[4]{\ensuremath{#1\,\models\,\set{#2}\;#3\;\set{#4}}}
\newcommand{\proofof}[1]{\ensuremath{\overline{#1}}}
\newcommand{\andif}{\ensuremath{\text{\;\;\;and\;\;\;}}}
\newcommand{\htriple}[3]{\ensuremath{\set{#1}\;#2\;\set{#3}}}
\newcommand{\triple}[3]{\ensuremath{\vdash\,\set{#1}\;#2\;\set{#3}}}
\newcommand{\truetriple}[3]{\ensuremath{\models\,\set{#1}\;#2\;\set{#3}}}

\newcommand{\safe}[2]{\ensuremath{\mathit{safe}_{#1}(#2)}}

\newcommand{\defined}[1]{\ensuremath{\mathsf{def}(#1)}}

\newcommand{\abort}{\ensuremath{\zorch}}

\newcommand{\live}[1]{\ensuremath{\mathit{live}(#1)}}

\newcommand{\femp}{\ensuremath{\mathbf{emp}}}
\newcommand{\fbar}[1]{\ensuremath{\mathbf{bar}_{#1}}}
\newcommand{\flock}[1]{\ensuremath{\mathbf{lock}_{#1}}}
\newcommand{\fexp}[1]{\ensuremath{\mathbf{exp}({#1})}}
\newcommand{\ftop}[1]{\ensuremath{\mathbf{top}_{#1}}}
%\newcommand{\fwrite}[1]{\ensuremath{\!\stackrel{#1}{\leadsto}\!}}
\newcommand{\fwrite}[1]{\ensuremath{\!\leadsto_{#1}\!}}
\newcommand{\fwrote}[1]{\ensuremath{\rightarrowtail_{#1}}}
\newcommand{\fpointsto}{\ensuremath{\!\mapsto\!}}
\newcommand{\fiter}[1]{\ensuremath{#1^{\ast}}}

\newcommand{\gensym}{\bullet}
\newcommand{\sepsym}{\ast}
\newcommand{\seqsym}{\vartriangleleft}
\newcommand{\sseqsym}{\blacktriangleleft}
\newcommand{\hashsym}{\#}

\DeclareMathOperator{\fhash}{\,\#\,}
\DeclareMathOperator{\fsep}{\,\ast\,}
\DeclareMathOperator{\fseq}{\,\vartriangleleft\,}
\DeclareMathOperator{\fsseq}{\,\blacktriangleleft\,}
\DeclareMathOperator{\fseqc}{\,\vartriangleright\,}
\DeclareMathOperator{\fsseqc}{\,\blacktriangleright\,}
\DeclareMathOperator{\fsepi}{\,\relbar\!\!\!\ast\,}
\DeclareMathOperator{\fseqli}{\,\relbar\!\vartriangleleft\,}
\DeclareMathOperator{\fseqri}{\,\vartriangleleft\!\relbar\,}
\DeclareMathOperator{\fsepic}{\,\ast\!\!\!\relbar\,}
\DeclareMathOperator{\fseqlic}{\,\vartriangleright\!\relbar\,}
\DeclareMathOperator{\fseqric}{\,\relbar\!\vartriangleright\,}
\DeclareMathOperator{\fgen}{\bullet}


\DeclareMathOperator{\shash}{\,\widetilde{\#}\,}
\DeclareMathOperator{\ssep}{\,\widetilde{\ast}\,}
\DeclareMathOperator{\sseq}{\,\widetilde{\vartriangleleft}\,}
\DeclareMathOperator{\ssseq}{\,\widetilde{\blacktriangleleft}\,}
\DeclareMathOperator{\sseqc}{\,\widetilde{\vartriangleright}\,}
\DeclareMathOperator{\ssseqc}{\,\widetilde{\blacktriangleright}\,}
\DeclareMathOperator{\ssepi}{\,\widetilde{\relbar\!\!\ast}\,}
\DeclareMathOperator{\sseqli}{\,\widetilde{\relbar\!\vartriangleleft}\,}
\DeclareMathOperator{\sseqri}{\,\widetilde{\vartriangleleft\!\relbar}\,}
\DeclareMathOperator{\ssepic}{\,\widetilde{\ast\!\!\relbar}\,}
\DeclareMathOperator{\sseqlic}{\,\widetilde{\vartriangleright\!\relbar}\,}
\DeclareMathOperator{\sseqric}{\,\widetilde{\relbar\!\vartriangleright}\,}
\DeclareMathOperator{\sgen}{\widetilde{\bullet}}

\DeclareMathOperator{\mhash}{\,\widehat{\#}\,}
\DeclareMathOperator{\msep}{\,\widehat{\ast}\,}
\DeclareMathOperator{\mseq}{\,\widehat{\vartriangleleft}\,}
\DeclareMathOperator{\msseq}{\,\widehat{\blacktriangleleft}\,}
\DeclareMathOperator{\mseqc}{\,\widehat{\vartriangleright}\,}
\DeclareMathOperator{\msseqc}{\,\widehat{\blacktriangleright}\,}
\DeclareMathOperator{\msepi}{\,\widehat{\relbar\!\!\ast}\,}
\DeclareMathOperator{\mseqli}{\,\widehat{\relbar\!\vartriangleleft}\,}
\DeclareMathOperator{\mseqri}{\,\widehat{\vartriangleleft\!\relbar}\,}
\DeclareMathOperator{\msepic}{\,\widehat{\ast\!\!\relbar}\,}
\DeclareMathOperator{\mseqlic}{\,\widehat{\vartriangleright\!\relbar}\,}
\DeclareMathOperator{\mseqric}{\,\widehat{\relbar\!\vartriangleright}\,}
\DeclareMathOperator{\mgen}{\widehat{\bullet}}


\newcommand{\amod}{\ensuremath{\mathcal{M}}}



\newcommand{\fbaru}{\ensuremath{\mathbf{bar}}}
\newcommand{\fwriteu}{\ensuremath{\!\leadsto\!}}

\newcommand{\pred}[1]{\ensuremath{[\![#1]\!]}}
\newcommand{\locked}[1]{\ensuremath{\mathit{locked}(#1)}}
\newcommand{\taurefines}{\ensuremath{\preceq}}
\newcommand{\modrefines}{\ensuremath{\leq}}
\newcommand{\lockcomplete}[1]{\ensuremath{\text{\textit{lock-complete}}(#1)}}
\newcommand{\stable}[1]{\ensuremath{\mathit{stable}(#1)}}
\newcommand{\stcong}[1]{\ensuremath{\sim_{#1}}}

%% Notation added for this report.
%% Notation that is not in the dissertation.

%% Rocq identifiers: typewriter, underscores allowed, may break at them.
\newcommand{\rocq}[1]{\nolinkurl{#1}}

%% A write by processor i that has flushed: (e ~>_i e') <| bar_i.
%% The dissertation's macro file already defines \fwrote{i} as a tailed
%% arrow with a subscript; we give it this meaning.
\newcommand{\fwroteu}{\ensuremath{\rightarrowtail}}

%% Ownership modes of the repaired judgment.
\newcommand{\modeov}{\ensuremath{\mathsf{ov}}}
\newcommand{\modedj}{\ensuremath{\mathsf{dj}}}
\newcommand{\mspec}[5]{\ensuremath{#1\,\vdash^{#2}\,\set{#3}\;#4\;\set{#5}}}
\newcommand{\mtruespec}[5]{\ensuremath{#1\,\models^{#2}\,\set{#3}\;#4\;\set{#5}}}

%% Rule names in running text.
\newcommand{\rn}[1]{\textsc{#1}}

%% Status notes while drafting. Remove before circulating.
\newcommand{\pending}[1]{\textcolor{red}{[\emph{pending:} #1]}}


\title{Weak-Memory Local Reasoning, Revisited}
\author{Ian Wehrman}
\date{Draft of \today}

\begin{document}
\maketitle

\begin{abstract}
My 2012 dissertation proposed a program logic for local reasoning about
concurrent programs on an x86-like, total-store-order memory model. It left
two components unfinished: the soundness proof, and a proof system for
assertions. Out of curiosity, I have returned to both with a mechanization
in Rocq. The soundness theorem turns out not to hold as stated, because four
of the original rules derive false specifications. However, a repaired
calculus is sound and adequate for the machine, and it keeps most of the
original design: the assertion language, the unit law for $\femp$, the
temporal frame rule, and the original load and store axioms. This report
recalls the original model and rules for a reader who knew the dissertation
but has forgotten its details. It then explains what was wrong, what had to
change and why, and what the repaired logic can prove.
\end{abstract}

\section{Introduction}
\label{sec:intro}

The dissertation~\cite{wehrman12} set out to give a local reasoning
principle for programs running on an x86-like machine, where each processor
buffers its writes and other processors cannot see a write until it leaves
the buffer. Its logic has three separating conjunctions. Spatial
conjunction $P \fsep Q$ splits memory by address, as in separation logic.
Temporal conjunction $P \fseq Q$ splits it by time: $P$ describes earlier
writes and $Q$ later ones, possibly to the same addresses. Strong temporal
conjunction $P \fsseq Q$ requires both. Each has a frame rule, and a
concurrency rule and a shared invariant in the style of concurrent
separation logic sit on top.

The dissertation was candid about where it stopped. Its statement of
contributions lists two major components as unfinished and admits a third
gap:
\begin{enumerate}[nosep]
  \item there is no proof system for the assertion language;
  \item the soundness proof for the specification logic is incomplete; and
  \item there are few worked examples, so it is unclear how one would
        actually use the rules to verify a program.
\end{enumerate}
It also hypothesized that the logic was sound and that the remaining proofs
were ``mostly straightforward inductions.''

I have wondered since whether that was true. This report returns to those
three items, fourteen years later, to satisfy that curiosity, with a
mechanization in the Rocq prover. Its findings are as follows.

\paragraph{The soundness theorem does not hold as stated.}
Several auxiliary claims are false, but each can be fixed without
changing the design (Section~\ref{sec:local-errors}). More seriously, the rules themselves are
inconsistent with any adequate semantics: from \rn{skip}, \rn{cons},
\rn{inv} and \rn{frame-tm} alone one can derive that a $\cskip$ changes the
contents of a memory cell (Section~\ref{sec:inconsistent}). No choice of
safety predicate can repair that.

\paragraph{The cause is general.}
An unrestricted temporal frame rule, a right unit for temporal conjunction,
the rule of consequence, and a command that gives away ownership of a
cell cannot coexist once some history that held before the transfer is
false after it (Section~\ref{sec:trilemma}). The observation is
independent of the memory model, and a sequential conjunction in which
later values override earlier ones supplies such histories.

\paragraph{A repaired calculus is sound.}
The repair replaces the invariant rule by lock and unlock rules that
transfer the invariant, restricts the temporal frame rule to commands that do not release the lock, and adjusts four side conditions
(Section~\ref{sec:changes}). The repaired calculus is proved sound by
induction on derivations, with an adequacy theorem for the dissertation's
machine (Section~\ref{sec:results}). The proofs are checked by the Rocq
kernel and use no axioms.

\paragraph{Most of the design survives.}
The machine, the assertion language and its interpretation of atoms, the
unit laws for $\femp$, overlapping temporal history, the temporal and
spatial frame rules, the concurrency rule, and the original load and store
axioms are all retained. So is the dissertation's assertion algebra,
with one family of exceptions: the original laws that rebracket or exchange
through strong temporal conjunction fail, and with them the strong temporal
frame rule. The exchange laws between spatial and temporal conjunction also
fail or are vacuous (Appendix~\ref{app:statements}).

\paragraph{Assertions have a proof system.}
A calculus of 47 rules for entailment between assertions is proved sound
(Section~\ref{sec:assertion-rules}). Its completeness is not established and was not
a goal, but it is enough: the larger derivations below use no semantic
entailment at all.

\paragraph{The logic can be used.}
Section~\ref{sec:proofs} derives a shared container in which a thread
acquires the lock, removes a cell, releases the lock, writes the cell
privately, and later returns it, while another processor's write stays
buffered throughout. It re-derives the dissertation's own message-passing
example. And it builds named locks on top of the single global lock, as
derived rules, and uses one to pass a message of any length from a producer
to a consumer.

\medskip
This report is for a reader who knew the original dissertation but has
forgotten its details. Section~\ref{sec:original} recalls the original model and
calculus in the original notation. Sections~\ref{sec:wrong}
and~\ref{sec:changes} say what was wrong and what changed.
Sections~\ref{sec:model} to~\ref{sec:proofs} present the repaired model,
calculus, results and derivations. Proofs are not reproduced; every claim
cites the name of its Rocq theorem, and the appendix gives fuller
statements.

\section{The original model and calculus}
\label{sec:original}

This section recalls Chapter~4 of the dissertation, in its notation. Nothing
here is new. Page and figure numbers refer to the dissertation.

\subsection{The machine}

A \emph{memory system} is a triple $\mu = (h,B,K)$: a heap
$h : \setlocations \pfun \setvalues$, an array of write buffers
$B : \setprocessors \tfun \listof{(\setlocations \times \setvalues)}$ with
the oldest write first, and a set $K \subseteq \setprocessors$ of
\emph{blocked} processors. A state pairs a stack $s$ with a memory system.
The lock is not a separate component: it is free when $K = \nil$ and held
by $i$ when $K = \setprocessors \setminus \set{i}$.

Primitive commands carry the identifier of the processor that executes
them. Their behaviour on processor $i$ is as follows.
\begin{itemize}[nosep]
  \item $\cstore{e}{e'}_i$ appends a write to $B(i)$. The address must be
        allocated, in the heap or in $B(i)$. A blocked processor may still
        store.
  \item $\cload{x}{e}_i$ requires $i \notin K$ and returns the newest write
        to that address in $B(i)$, or else the heap value.
  \item $\cfence_i$, $\clock_i$ and $\cunlock_i$ require $B(i) = \lnil$.
        $\clock_i$ additionally requires $K = \nil$ and sets
        $K = \setprocessors \setminus \set{i}$; $\cunlock_i$ reverses this.
  \item A silent step $\taustep$ moves the oldest write of some unblocked
        processor from its buffer to the heap.
\end{itemize}
Commands are built from primitives by sequencing $\cseq{c}{c'}$, choice
$\cchoice{c}{c'}$, iteration $\cloop{c}$ and parallel composition
$\cpar{c}{c'}$. There is no allocation. While a processor holds the lock,
no other processor can load or flush, but others may continue to store and
to assign to registers.

\subsection{Generalized memories and separation}

A \emph{generalized memory system} $\nu = (h,B,K,\Gamma)$ adds a
\emph{buffer-completeness} set $\Gamma \subseteq \setprocessors$, with
$h = \nil$ whenever $\Gamma = \nil$. Membership $i \in \Gamma$ means that
buffer $i$ is complete: nothing may be placed before it.

Three operations combine generalized memories. Write $\nu_k =
(h_k,B_k,K_k,\Gamma_k)$. All three take $K_0 \cup K_1$ and
$\Gamma_0 \cup \Gamma_1$ and require $K_0 \cap K_1 = \nil$.
\begin{description}[nosep,leftmargin=1.5em]
  \item[Spatial, $\nu_0 \msep \nu_1$.] The address domains of the two
    sides, counting both heaps and buffers, are disjoint. The heap is
    $h_0 \uplus h_1$ and each buffer is an interleaving of the two.
  \item[Temporal, $\nu_0 \mseq \nu_1$.] The heap is $h_0 \override h_1$, so
    later values win, and each buffer is $B_0(i) \lapp B_1(i)$. The sides
    may share addresses, subject to a compatibility condition: an unblocked
    pending write in $\nu_0$ may not target an address in the heap of
    $\nu_1$, nor an address pending on a different unblocked processor in
    $\nu_1$. Finally $B_0(i) = \lnil$ for every $i \in \Gamma_1$.
  \item[Spatiotemporal, $\nu_0 \mhash \nu_1$.] As temporal, but each buffer
    is any interleaving whose lookup function agrees with
    $B_0(i) \lapp B_1(i)$, and the condition on $\Gamma_1$ is dropped. It
    contains both of the others and is used only in the semantics of
    specifications.
\end{description}
\emph{Strong temporal} separation $\nu_0 \msseq \nu_1$ is $\nu_0 \mseq
\nu_1$ restricted to pairs that are also spatially compatible.

\subsection{Assertions}

Assertions are given by
\begin{align*}
  P \bnfdef{} & b \bnfbar P \disj P \bnfbar P \conj P \bnfbar \exists x \st P
    \bnfbar \forall x \st P \bnfbar \femp \bnfbar \fbar{e} \bnfbar
    \flock{e} \bnfbar e \fwrite{e'} e'' \\
  \bnfbar{} & P \fsep P \bnfbar P \fseq P \bnfbar P \fsepi P \bnfbar
    P \fseqli P \bnfbar P \fseqri P
\end{align*}
and are interpreted over a stack and a generalized memory. The atoms are:
\begin{itemize}[nosep]
  \item $\femp$: empty heap, empty buffers, $K = \nil$; $\Gamma$ is
        unconstrained.
  \item $\fbar{i}$: as $\femp$, with $i \in \Gamma$ when $i$ is a processor.
  \item $\flock{i}$: empty heap and buffers, with
        $K = \setprocessors \setminus \set{i}$.
  \item $e \fwrite{i} e'$ (``leads to''): either a single \emph{pending}
        write $(e,e')$ in $B(i)$ with an empty heap and
        $K \subseteq \set{i}$, or a single \emph{flushed} heap cell
        $e \mapsto e'$ with empty buffers, $K = \nil$ and $i \in \Gamma$.
\end{itemize}
$P \fsep Q$ and $P \fseq Q$ hold of $\nu$ when $\nu$ is a spatial,
respectively temporal, combination of a model of $P$ and a model of $Q$.
The three implications are the corresponding adjoints. Two abbreviations
matter below:
\[
  P \fsseq Q \eqdef (P \fsep Q) \conj (P \fseq Q)
  \qquad\qquad
  e \fpointsto e' \eqdef (\exists x \st e \fwrite{x} e') \fseq
                         (\forall x \st \fbar{x}).
\]
Points-to therefore describes a flushed cell with \emph{every} processor
buffer-complete. The dissertation claims that every assertion is closed
under flushing (Proposition~9), and that each conjunction is associative
with $\femp$ as a unit (Figure~4.5).

\subsection{Specifications and rules}

A specification $\spec{J}{P}{c}{Q}$ has an invariant $J$ for shared state
and a private precondition and postcondition. Figure~\ref{fig:original}
reproduces the axioms and rules.

\begin{figure}[t]
\small
\begin{mathpar}
\inferrule[skip]{ }{\spec{J}{P}{\cskip_i}{P}}
\and
\inferrule[assume]{ }{\spec{J}{\neg b \disj P}{\cassume{b}_i}{P}}
\and
\inferrule[assert]{ }{\spec{J}{b \conj P}{\cassert{b}_i}{P}}
\and
\inferrule[assign]{ }{\spec{J}{P\subst{e}{x}}{\cassign{x}{e}_i}{P}}
\and
\inferrule[fence]{ }{\spec{J}{\femp}{\cfence_i}{\fbar{i}}}
\and
\inferrule[lock]{ }{\spec{J}{\femp}{\clock_i}{\flock{i}}}
\and
\inferrule[unlock]{ }{\spec{J}{\flock{i}}{\cunlock_i}{\femp}}
\and
\inferrule[load]{ }
  {\spec{J}{e \fwrite{i} e' \fsseq P}{\cload{x}{e}_i}
        {(e \fwrite{i} e' \fsseq P) \conj x = e'}}
\and
\inferrule[store]{ }
  {\spec{J}{e \fwrite{i} e'' \fsseq P}{\cstore{e}{e'}_i}
        {(e \fwrite{i} e'' \fsseq P) \fseq e \fwrite{i} e'}}
\and
\inferrule[frame-$\bullet$]
  {\spec{J}{P}{c}{Q} \\ \mod{c} \cap \fv{R} = \nil}
  {\spec{J}{R \fgen P}{c}{R \fgen Q}}
\and
\inferrule[cons]
  {P \sentails P' \\ \spec{J}{P'}{c}{Q'} \\ Q' \sentails Q}
  {\spec{J}{P}{c}{Q}}
\and
\inferrule[seq]
  {\spec{J}{P}{c}{R} \\ \spec{J}{R}{c'}{Q}}
  {\spec{J}{P}{\cseq{c}{c'}}{Q}}
\and
\inferrule[choice]
  {\spec{J}{P}{c}{Q} \\ \spec{J}{P}{c'}{Q}}
  {\spec{J}{P}{\cchoice{c}{c'}}{Q}}
\and
\inferrule[loop]{\spec{J}{P}{c}{P}}{\spec{J}{P}{\cloop{c}}{P}}
\and
\inferrule[disj]
  {\spec{J}{P}{c}{Q} \\ \spec{J}{P'}{c}{Q}}
  {\spec{J}{P \disj P'}{c}{Q}}
\and
\inferrule[ex]
  {\spec{J}{P}{c}{Q} \\ x \notin \fv{c,Q}}
  {\spec{J}{\exists x \st P}{c}{Q}}
\and
\inferrule[conc]
  {\spec{J}{P_k}{c_k}{Q_k} \\ \fv{P_k,c_k,Q_k} \cap \mod{c_{1-k}} = \nil
   \quad (k = 0,1)}
  {\spec{J}{P_0 \fsep P_1}{\cpar{c_0}{c_1}}{Q_0 \fsep Q_1}}
\and
\inferrule[share]
  {\spec{J}{P}{c}{Q}}
  {\spec{\femp}{J \fsep P}{c}{J \fsep Q}}
\and
\inferrule[inv]
  {\spec{\femp}{J \fsep P \fsep \flock{i}}{c}{J \fsep Q \fsep \flock{i}}}
  {\spec{J}{P \fsep \flock{i}}{c}{Q \fsep \flock{i}}}
\end{mathpar}
\caption{The original axioms and rules (Figures~4.7 and~4.8).
\rn{frame-$\bullet$} stands for three rules, \rn{frame-sp}, \rn{frame-tm}
and \rn{frame-stm}, with $\bullet$ read as $\fsep$, $\fseq$ and $\fsseq$.}
\label{fig:original}
\end{figure}

Three features of the design are worth recalling, because the repair turns
on them. First, \rn{lock} and \rn{unlock} produce and consume only the
token $\flock{i}$. The invariant reaches the lock holder through \rn{inv},
which may be applied to any command executed while the lock is held and
allows the private state to exchange cells with $J$. Second, \rn{frame-tm}
has only the usual variable condition: any history $R$ may be placed before
any specification. Third, temporal conjunction is not cancellative, as the
dissertation remarks: $\mu \sseq \mu_0 = \mu \sseq \mu_1$ does not imply
$\mu_0 = \mu_1$.

\subsection{Semantics and status}

Truth of a specification is defined through a family of predicates
$\safe{n}{c,s,\mu,J,Q}$ after Vafeiadis. The predicate holds at $n+1$ when:
if $c$ is $\cskip$ then $(s,\mu,\setprocessors) \sentails Q$; and for every
$\mu_J$, $\mu_F$ and $\mu_0 \in \mu_J \ssep (\mu_F \shash \mu)$ that is
lock-complete, with either $(s,\mu_J,\setprocessors) \sentails J$ or the
lock held in $\mu_0$, the configuration does not abort, and every step to
$(c',s',\mu_1)$ admits a new decomposition
$\mu_1 \in \mu'_J \ssep (\mu'_F \shash \mu')$ with
$\mu'_F \taurefines \mu_F$, the same condition on $J$, and
$\safe{n}{c',s',\mu',J,Q}$. A specification is true,
$\truespec{J}{P}{c}{Q}$, when $(s,\mu,\setprocessors) \sentails P$ implies
$\safe{n}{c,s,\mu,J,Q}$ for all $n$. Local states are always read at full
completeness, $\Gamma = \setprocessors$.

Theorem~1 states that provable specifications are true. The dissertation
gives soundness lemmas for the axioms in an appendix, sketches a frame lemma for the
spatiotemporal operator $\shash$, and reduces the spatial frame rule to Conjecture~1, which is left open.
The conjecture says that spatial separation survives a step: if a memory
splits spatially into $\mu_b$ and $\mu_c$ and takes a step, then any
spatiotemporal split of the result into a flushed descendant of $\mu_b$ and
some $\mu'_c$ is again a spatial split. Associativity of buffer overriding
(Proposition~2) and of the three separations, and flushing closure of
assertions (Proposition~9), are stated with at most a sketch.

\section{What was wrong}
\label{sec:wrong}

The problems are presented from the smallest to the most serious,
which is also the order in which the mechanization found them. Each is a
theorem checked by Rocq, and the theorem's name is given so that the
statement can be inspected. All of them are stated against a transcription of the original definitions.

\subsection{Errors that can be fixed in place}
\label{sec:local-errors}

Table~\ref{tab:local} lists six original claims that are false as stated.
None is deep: each is fixed by correcting the one definition or side
condition involved, and the rest of the design is unaffected.



\begin{table}[!ht]
\small
\begin{tabularx}{\textwidth}{@{}p{0.27\textwidth}X p{0.23\textwidth}@{}}
\toprule
Original claim & Counterexample & Repair \\
\midrule
\rn{assign} is sound for every invariant
  & With $J$ the assertion $x = 0$, the axiom derives
    $\spec{J}{\bexpt}{\cassign{x}{1}_1}{\bexpt}$, and the assignment
    falsifies $J$ with the lock free.
    \rocq{printed_assignment_invalid}
  & Require $x \notin \fv{J}$. \\[2pt]
Buffer overriding is associative (Prop.~2)
  & Three one-element buffers writing $1$, $2$, $3$ to one address: the
    list $[2,1,3]$ is in $a \override (b \override c)$ but not in
    $(a \override b) \override c$.
    \rocq{printed_buffer_override_not_associative}
  & A write from the later buffer may be taken only after the earlier
    buffer's writes to that address. \\[2pt]
Spatiotemporal separation is associative
  & Compatibility ignores blocked processors, so whether two memories
    compose can depend on a third that supplies the blocking fact.
    \rocq{printed_spatiotemporal_not_associative}
  & Check conflicts whether or not the processor is blocked. \\[2pt]
Assertions are closed under flushing (Prop.~9)
  & $\bexpf \fseqri \fbar{1}$ holds vacuously while processor $1$ has a
    pending write, and fails once the write flushes.
    \rocq{original_flushing_closure_false}
  & Implications quantify over all futures of the current memory. \\[2pt]
Conjecture~1
  & A cell duplicated across $\shash$ with equal values is not a spatial
    split. \rocq{printed_conjecture1_false}
  & Not needed by the final proof. \\[2pt]
\rn{unlock} is sound for every invariant
  & With $J = \bexpf$ and the lock held, $\cunlock_i$ is enabled and
    nothing restores $J$. \rocq{original_arbitrary_invariant_unlock_invalid}
  & See Section~\ref{sec:changes}. \\
\bottomrule
\end{tabularx}
\caption{Original claims that fail, each with a checked counterexample.}
\label{tab:local}
\end{table}

Two remarks. The repaired compatibility relation is stronger than the
original one, so it rejects some compositions the dissertation admits. Those
are states in which a blocked processor has a pending write to a cell that
someone else owns. And the repaired implications are
the ordinary Kripke reading of implication over the flushing preorder; the
three adjunctions between conjunctions and implications still hold.

The failure of \rn{unlock} points at something larger. The original safety
predicate asks for the invariant only while the lock is free. Once a
thread holds the lock, the predicate has forgotten that $J$ held at the
moment of acquisition, so neither \rn{unlock} nor \rn{inv} can be justified
from it; \rn{inv} on its own takes a true specification to a false one
(\rocq{original_inv_alone_counterexample}). That alone would call for a different predicate. The next
subsection shows that it calls for different rules.

\subsection{The original rules are inconsistent}
\label{sec:inconsistent}

The derivations below use the dissertation's own points-to, $x \fpointsto
v$, and are checked against its own interpretation of assertions.

\paragraph{History can be invented.}
A model is a memory and records nothing of how that memory came about. So a
current cell is also a model of any stale value followed by that cell:
\begin{equation}
  \label{eq:invent}
  x \fpointsto 7 \fsep \flock{i}
  \;\sentails\;
  x \fpointsto 5 \fseq (x \fpointsto 7 \fsep \flock{i}).
\end{equation}
The stale cell is overridden by the later one, and nothing remains of it.
This is the non-cancellativity that the dissertation notes in passing.

\paragraph{A $\cskip$ that changes memory.}
Let $J$ be $\femp \disj x \fpointsto 7$. The following is a derivation in
the original calculus.
\begin{enumerate}[nosep]
  \item $\spec{\femp}{J \fsep x \fpointsto 7 \fsep \flock{i}}{\cskip_i}
        {J \fsep \femp \fsep \flock{i}}$,
        by \rn{skip} and \rn{cons}: $J$ cannot be its second disjunct, since
        the cell would occur twice, and in the first disjunct the cell can
        be reassigned to $J$.
  \item $\spec{J}{x \fpointsto 7 \fsep \flock{i}}{\cskip_i}
        {\femp \fsep \flock{i}}$, by \rn{inv}. The thread has given its
        cell to the invariant.
  \item $\spec{J}{x \fpointsto 5 \fseq (x \fpointsto 7 \fsep \flock{i})}
        {\cskip_i}{x \fpointsto 5 \fseq (\femp \fsep \flock{i})}$,
        by \rn{frame-tm}.
  \item $\spec{J}{x \fpointsto 7 \fsep \flock{i}}{\cskip_i}
        {x \fpointsto 5 \fsep \flock{i}}$, by \rn{cons}, using
        \eqref{eq:invent} on the left and the unit law for $\femp$ on the
        right.
\end{enumerate}
The conclusion says that a thread holding the lock and a cell containing
$7$ may do nothing and then own that cell containing $5$. In the state
where processor $1$ holds the lock and $x$ contains $7$, the precondition
holds and the postcondition does not.
(\rocq{printed_points_stale_value_unsound}; the mechanized statement parses
\rn{inv}'s premise as $J \fsep (P \fsep \flock{i})$.)

\paragraph{A $\cskip$ that drains another processor's buffer.}
Spatial conjunction takes the union of the completeness sets and does not
require them to be disjoint, so an empty neighbour can copy a cell's marks:
$x \fpointsto 7 \fsep \femp \sentails x \fpointsto 7 \fsep \fbar{2}$. With
$J$ the cell $x \fpointsto 7$, \rn{skip}, \rn{cons} and \rn{inv} give
$\spec{J}{\femp \fsep \flock{1}}{\cskip_1}{\fbar{2} \fsep \flock{1}}$, and
\rn{frame-tm} with a pending write $R = y \fwrite{2} 17$ gives
\[
  \spec{J}{R \fseq (\femp \fsep \flock{1})}{\cskip_1}
          {R \fseq (\fbar{2} \fsep \flock{1})}.
\]
When processor $1$ holds the lock and processor $2$'s write to $y$ is still
buffered, the precondition holds. The postcondition places a barrier for
processor $2$ after its pending write, which that memory does not satisfy.
(\rocq{printed_points_foreign_barrier_unsound}.)

\paragraph{The strong temporal frame rule does not escape.}
One might hope that \rn{frame-stm} is safe, since its frame is also
spatially separate. It is not, because $R \fsseq P$ abbreviates
$(R \fsep P) \conj (R \fseq P)$ and the two conjuncts may be witnessed by
different splits. Take $R = \femp \disj x \fpointsto 5$. In the state above,
the spatial conjunct of $R \fsseq (x \fpointsto 7 \fsep \flock{1})$ holds
with $R$ empty, and the temporal conjunct holds with $R$ the stale cell.
Applying \rn{frame-stm} to step~2 yields a postcondition that the unchanged
state cannot satisfy. (\rocq{printed_strong_frame_counterexample}; this
derivation uses \rn{inv} and \rn{frame-stm} but not \rn{frame-tm}.)

All three conclusions are also invalid for the original safety predicate,
but that is the lesser point. They are false under the
dissertation's informal reading of a specification, and under any semantics
for which a true specification describes what the program does.

\subsection{Why no semantics could repair it}
\label{sec:trilemma}

The first derivation is an instance of a general conflict.

\begin{proposition}[\rocq{trilemma}]
\label{prop:trilemma}
Suppose a logic of triples has an unrestricted temporal frame rule, the
rule of consequence, and the entailment $R \fseq \femp \sentails R$.
Suppose $\htriple{P}{c}{\femp}$ is provable. Then $\htriple{R \fseq P}{c}{R}$
is provable for every $R$. Hence if some execution of $c$ starts in a state
satisfying $R \fseq P$ and ends in a state that does not satisfy $R$, the
logic is unsound.
\end{proposition}

The proof is two rule applications. The content is in the last sentence. A
triple $\htriple{P}{c}{\femp}$ with $P$ owning a cell is a \emph{transfer-out
contract}: after $c$, the cell belongs to someone else. Whenever temporal
conjunction lets later values override earlier ones, $R$ may describe an
older state of the same cell, as the store axiom's own postcondition does.
Framing $R$ across the transfer then hands the thread a stale copy of a
cell it no longer owns.

So four things cannot all be kept: the unrestricted temporal frame rule,
the right-unit law, transfer-out contracts ending in $\femp$, and
overlapping temporal history. A sound logic must give one up, and the
choice determines its character:
\begin{itemize}[nosep]
  \item \emph{No transfer out.} Cells may be taken from the invariant and
        never returned. This is sound and useless for a shared container.
  \item \emph{No right unit.} Let $\femp$ record that a cell was given
        away. Framing survives, and every postcondition that mentions
        $\femp$ changes meaning.
  \item \emph{Transfer leaves a mark.} Return an explicit tombstone
        assertion instead of $\femp$. The algebra survives; every later
        specification carries the tombstones.
  \item \emph{Restrict the frame.} Forbid temporal frames that overlap the
        cells being given away. Everything else is kept.
\end{itemize}
The first two were built and abandoned during this work, for the reasons
given. The repaired calculus takes the fourth.
Proposition~\ref{prop:trilemma} shows that some restriction of this kind is
needed, and the derivations above supply histories that do become false
after a transfer. It does not show that the particular restriction adopted
is the weakest possible.

The second derivation has a different cause and a different repair. The
all-processor completeness of points-to, and more generally the
completeness marks of whatever the invariant holds, should not pass to a
thread that acquires the lock. A thread can vouch only for its own buffer.

\subsection{The dissertation's own example}
\label{sec:own-example}

Section~4.1 of the dissertation proves a small message-passing program
correct, using an invariant that mentions two registers and the derived
rule \rn{atomic}. That proof does not carry over to the repaired calculus:
\rn{atomic} and \rn{share} in their original forms, and
auxiliary-variable elimination, are not available, and the concurrency rule now forbids either thread from
modifying a variable free in the invariant. The proof does not carry
over, although its conclusion is still true.
Section~\ref{sec:own-example-repaired} derives the same program from a
different invariant and recovers the original conclusion.

\section{What changed}
\label{sec:changes}

\pending{to be written}

\section{The repaired model}
\label{sec:model}

\pending{to be written}

\section{The repaired calculus}
\label{sec:calculus}

\pending{to be written}

\section{Soundness and adequacy}
\label{sec:results}

\pending{to be written}

\section{How proofs proceed}
\label{sec:proofs}

The dissertation gave few examples. This section gives three, in increasing
size. All are derivations in the calculus of Section~\ref{sec:calculus},
checked by Rocq.

\subsection{A store, with and without a fence}

For any invariant and either mode,
\[
  \mspec{J}{m}{x \fcell u}{\cseq{\cstore{x}{v}_i}{\cfence_i}}{x \fcell v}
\]
is derivable (\rocq{owned_derive_store_fence_cell}): \rn{store-any} gives
$x \fcell u \fseq x \fwrite{i} v$, \rn{frame-tm} carries that history
across $\cfence_i$, and a drained history collapses to its last write. This
is the store rule of ordinary separation logic. The same specification
without the fence is \emph{invalid} (\rocq{owned_plain_store_cell_invalid}):
the store may terminate with its write still buffered, and a cell assertion
has no buffered writes. The logic is the difference between these two.

\subsection{Reading a buffered write}

After $\cstore{x}{v}_i$ the thread holds $R \fseq x \fwrite{i} v$, where
$R$ describes the older state of $x$. The load axiom wants the loaded write
to be spatially separate from the rest, and an old cell cannot be separated
spatially from a newer write to the same address. The temporal frame rule
is what resolves this. The newest write alone satisfies the axiom's
precondition, with $\femp$ for $P$; the older history is a temporal frame:
\[
\inferrule*[right=frame-tm]
  {\mspec{J}{\modeov}{x \fwrite{i} v}{\cload{r}{x}_i}
         {x \fwrite{i} v \conj r = v}}
  {\mspec{J}{\modeov}{R \fseq x \fwrite{i} v}{\cload{r}{x}_i}
         {R \fseq (x \fwrite{i} v \conj r = v)}}
\]
(\rocq{owned_derive_load_after_history}). Store-to-load forwarding, the
behaviour most characteristic of the machine, is thus an instance of the
rule that most distinguishes the logic, and the reason it was worth
keeping.

\subsection{A shared container}

One cell $x$ is shared through a flag $f$. With
$X = \exists v \st x \fcell v$, the invariant is
\[
  J \;=\; f \fcell 0 \;\disj\; (f \fcell 1 \fsep X).
\]
A thread on processor $1$ takes the item, writes to it privately, and
returns it:
\begin{align*}
 & \clock_1;\ \cload{r}{f}_1;\ \cassume{r = 1}_1;\ \cstore{f}{0}_1;\
   \cunlock_1; \\
 & \cstore{x}{9}_1; \\
 & \clock_1;\ \cload{r}{f}_1;\ \cassume{r = 0}_1;\ \cstore{f}{1}_1;\
   \cunlock_1 .
\end{align*}
Meanwhile processor $2$ has an unrelated write $R = y \fwrite{2} 17$
buffered. The derived specification is
$\mspec{J}{\modedj}{R}{c}{R}$ (\rocq{owned_container_derived}). Writing
$H(a,v)$ for $(\exists u \st a \fcell u) \fseq a \fwrite{1} v$, the proof
passes through these assertions, each with $R$ as a spatial frame:
\begin{center}\small
\begin{tabular}{@{}ll@{}}
\toprule
After & The thread holds \\
\midrule
first load and assume & $\flock{1} \fsep (f \fcell 1 \fsep X)$ \\
$\cstore{f}{0}$ & $\flock{1} \fsep (H(f,0) \fsep X)$ \\
first unlock & $X$ \\
$\cstore{x}{9}$ & $H(x,9)$ \\
second lock & $H(x,9) \fseq ((\flock{1} \fsep J) \fseq \fbar{1})$ \\
second load and assume & $\flock{1} \fsep (f \fcell 0 \fsep x \fcell 9)$ \\
$\cstore{f}{1}$ & $\flock{1} \fsep (H(f,1) \fsep x \fcell 9)$ \\
second unlock & $\femp$ \\
\bottomrule
\end{tabular}
\end{center}
Three things are visible in it.

\emph{The foreign write costs one rule application.} $R$ concerns an
address the thread never touches, so a single \rn{frame-sp} around the
entire program carries it, across both releases.

\emph{The temporal frame rule carries the private write across the second
acquisition.} $H(x,9)$ is overlapping history, possibly still buffered.
\rn{frame-tm} places it before the acquired state, and the barrier that
\rn{lock} delivers then drains it, leaving $x \fcell 9$.

\emph{Releases use entailment, not framing.} The two \rn{unlock} premises
are $(H(f,0) \fsep X) \fseq \fbar{1} \sentails J \fsep X$ and
$(H(f,1) \fsep x \fcell 9) \fseq \fbar{1} \sentails J \fsep \femp$. The
first keeps the item and the second donates it.

The example is not vacuous. From a memory with the lock free, $f = 1$,
$x = 7$ and processor $2$'s write pending, there is a complete 24-step
execution in which that write is pending in every state
(\rocq{owned_container_foreign_pending_throughout}), and by adequacy every
execution from that memory is safe and, if it completes, restores the
invariant.

A second version replaces each $\cassume{}$ by a loop that releases the
lock and retries when the flag is wrong (\rocq{owned_retrying_container_derived}).
Its return loop has a branch for ``the container is already full'', and
that branch is derived from $\bexpf$: the thread holds $x$, so the pool
cannot.

\paragraph{The cost.}
The container's derivation applies \rn{cons} ten times and needs sixteen
entailments that are not identities, proved here semantically. Most are
rearrangements, a few are the laws about barriers and histories used
above, and one, at the second acquisition, takes fifty lines.
\pending{the same derivation with syntactic entailments; the Section 4.1
example; named locks and a producer/consumer client}

\section{The mechanization}
\label{sec:mechanization}

\pending{to be written}

\section{Later work}
\label{sec:related}

This is a brief orientation, and it relies on the papers' own accounts of
their results.

\paragraph{Before and alongside.}
Ridge~\cite{ridge10} gave a rely-guarantee proof system for x86-TSO,
mechanized in HOL4 and applied to Simpson's four-slot algorithm. Cohen and
Schirmer~\cite{cohen10} and Owens~\cite{owens10} took the other route:
conditions on a program, an ownership discipline and triangular-race
freedom respectively, under which its TSO executions are sequentially
consistent, so that buffers need not be reasoned about at all. The
dissertation, and the proposal with Berdine that preceded
it~\cite{wehrman11}, sat between these: buffers are visible in assertions,
and separation is meant to keep the reasoning local.

\paragraph{iCAP-TSO.}
The closest successor is the logic of Sieczkowski, Svendsen, Birkedal and
Pichon-Pharabod~\cite{sieczkowski15}, which cites the dissertation. It
extends a higher-order concurrent separation logic with store buffers and
with connectives for them, including one stating that an assertion holds
until a buffered update is flushed and another afterwards. That is a
temporal connective in the dissertation's sense, attached to a single
update. Their central idea is a \emph{fiction of sequential consistency}:
a library that synchronizes enough can be given a specification from which
its clients reason in ordinary separation logic. They observe, fairly, that
the dissertation's logic cannot exploit synchronization to simplify client
proofs in that way. They also verify a spin lock whose release is a plain
buffered store, which is exactly the unlocked access to shared state that
Section~\ref{sec:open} leaves open here. Jacobs~\cite{jacobs14} proposed
``TSO spaces'' for separation logic with a related aim.

\paragraph{Temporal reasoning elsewhere.}
Bornat, Alglave and Parkinson~\cite{bornat15} built a rely-guarantee logic
for weak memory around temporal modalities, ``since'' and ``so far'' among
them. The idea that ordering between writes belongs in the assertion
language, which the dissertation expressed as a separating conjunction, has
thus reappeared in other forms.

\paragraph{Other memory models.}
Most later work targets the C11 model and weaker ones rather than TSO:
relaxed separation logic~\cite{vafeiadis13} and GPS~\cite{turon14}, and
their reconstruction inside Iris~\cite{kaiser17}, which made mechanized
soundness proofs the norm for such logics. For x86 itself,
Pierogi~\cite{bila22} is an Owicki--Gries logic over a view-based semantics
that also covers persistent memory.

\paragraph{Where this report stands.}
Against that background, what the repaired logic offers is narrow and, as
far as I know, still distinctive: sequential and spatial separating
conjunctions over the whole history of a thread's writes, each with a frame
rule, and a soundness proof checked end to end against an operational TSO
machine. Section~\ref{sec:trilemma}'s conflict between temporal framing,
units and ownership transfer does not depend on TSO and may be of use to
anyone designing a logic with a sequential conjunction. What it lacks,
unlocked access and a way for libraries to hide the memory model from
their clients, is what the later logics were built to provide.

\section{What remains open}
\label{sec:open}

\pending{to be written}


\begin{thebibliography}{99}
\bibitem{bila22}
E.~V. Bila, B.~Dongol, O.~Lahav, A.~Raad and J.~Wickerson.
\newblock View-based Owicki--Gries reasoning for persistent x86-TSO.
\newblock In \emph{ESOP}, 2022.
\bibitem{bornat15}
R.~Bornat, J.~Alglave and M.~Parkinson.
\newblock New lace and arsenic: adventures in weak memory with a program
  logic.
\newblock arXiv:1512.01416, 2015.
\bibitem{cohen10}
E.~Cohen and B.~Schirmer.
\newblock From total store order to sequential consistency: a practical
  reduction theorem.
\newblock In \emph{ITP}, 2010.
\bibitem{jacobs14}
B.~Jacobs.
\newblock Verifying TSO programs.
\newblock Report CW660, KU Leuven, 2014.
\bibitem{kaiser17}
J.-O. Kaiser, H.-H. Dang, D.~Dreyer, O.~Lahav and V.~Vafeiadis.
\newblock Strong logic for weak memory: reasoning about release-acquire
  consistency in Iris.
\newblock In \emph{ECOOP}, 2017.
\bibitem{owens10}
S.~Owens.
\newblock Reasoning about the implementation of concurrency abstractions on
  x86-TSO.
\newblock In \emph{ECOOP}, 2010.
\bibitem{ridge10}
T.~Ridge.
\newblock A rely-guarantee proof system for x86-TSO.
\newblock In \emph{VSTTE}, 2010.
\bibitem{sieczkowski15}
F.~Sieczkowski, K.~Svendsen, L.~Birkedal and J.~Pichon-Pharabod.
\newblock A separation logic for fictional sequential consistency.
\newblock In \emph{ESOP}, 2015.
\bibitem{turon14}
A.~Turon, V.~Vafeiadis and D.~Dreyer.
\newblock GPS: navigating weak memory with ghosts, protocols, and
  separation.
\newblock In \emph{OOPSLA}, 2014.
\bibitem{vafeiadis13}
V.~Vafeiadis and C.~Narayan.
\newblock Relaxed separation logic: a program logic for C11 concurrency.
\newblock In \emph{OOPSLA}, 2013.
\bibitem{wehrman11}
I.~Wehrman and J.~Berdine.
\newblock A proposal for weak-memory local reasoning.
\newblock In \emph{LOLA}, 2011.
\bibitem{wehrman12}
I.~Wehrman.
\newblock \emph{Weak-Memory Local Reasoning}.
\newblock PhD dissertation, The University of Texas at Austin, December 2012.
\end{thebibliography}

\appendix
\section{Statements in full}
\label{app:statements}

\pending{fuller statements of the principal definitions and theorems, with
their Rocq names; the table of the dissertation's assertion laws}


\end{document}
